\chapter{存在与坍塌：当作用量趋近于零}

位移、制造和遗传都要依赖实际条件。物理稳定、功能维持与社会责任采用不同判据：一块岩石可以被动稳定，一种服务需要持续运行，一个法律关系需要制度支持。技术改变其中一段的成本，会重新安排关联结构，也让“它以什么方式存在、凭什么被依赖”成为具体问题。

生命之所以会对干预作出反应，是因为它必须持续维持自身。这种反应不是被动受力，而是把外部作用转成自身维持过程的一部分。把这一点放到更长的时间尺度上，它并不只是生物学事实。任何能够持续存在的东西，都在以某种方式对抗溃散，而维持自身的代价由什么承付，决定了它在现实中占有多重的分量。

\section{存在的判据：什么算存在，什么可以被依赖}

存在、知识与责任分别要求不同的判据。一个对象可以真实存在，却不构成可靠证据；一个判断可以成立，却没有获得执行授权；一个系统可以造成后果，却没有相应的承受边界。机器生成把这几种关系推到前台。区分它们，才能说明轻载体如何抵达重后果，以及哪一步缺少了必要的凭据或责任。

\subsection{存在的重量：从物理稳定到智能生成}

\subhead{稳定与切分}

世界最初没有图像，也没有名字，甚至没有我们今天意义上的“物”。在极早的宇宙中，结构尚未稳定，能量、场、粒子和涨落交织在一起，存在还不是对象，而是连续的相互作用和不稳定的差异。后来，宇宙膨胀，温度下降，原子形成，恒星点燃，元素合成，行星冷却，世界才逐渐从混沌的能量状态中沉淀出稳定结构。存在的第一步，不是被命名，而是稳定下来；没有稳定结构，就没有可以被反复指认的东西，也没有“这个东西”与“那个东西”的区别。

稳定有不同的维持条件。岩石的结构取决于化学键、晶格和环境，行星轨道取决于运动状态与引力关系；这些稳定未必要求持续的外部做功。生命则通过代谢、修复和调节维持远离平衡的状态，组织依靠人员、记录与制度持续运行。存在的物理条件、功能的维持成本和社会责任由此分开，后面要分别追问它们如何被支持。

生命出现以后，世界不再只是外部环境。最早的生命并不拥有世界模型，它只是在化学梯度中维持边界，摄入能量，排出废物，修复损伤，趋近有利条件，远离破坏条件。对生命来说，感知不是为了描绘世界，而是为了继续存在。同一片水域，对岩石只是环境；对生命体，却可能是资源、威胁、温度、压力和通道。世界没有改变，但某些差异开始被生命内部的维持过程选出来。不是所有存在都同样重要，只有那些影响代谢、修复、移动和生存的差异，才进入生命的世界。

智能的本体论根源就在这里。智能不是凭空面对一个已经切分好的对象世界；现实本身是连续的过程、流动、扰动、能量差和物质分布。智能首先要做的，不是求解，而是在连续现实中切出可行动的差异：什么可以接近，什么必须避开；什么可以摄入，什么会伤害；什么是路径，什么是障碍；什么是自身，什么是外部。所谓“对象”，不是世界天然贴好的标签，而是稳定结构与行动需求相遇后形成的切分。生命按代谢和危险切分，动物按路径、猎物和威胁切分，人类按名称、工具、制度和责任切分。智能的差异，首先不是答案不同，而是把什么东西当成“一个东西”的方式不同。

表示由此萌芽。足迹、气味、巢穴、叫声、姿态和迁徙路线，不是图像，也不是文字，却已经让不在场的东西继续发生作用。一个气味标记让离开的身体仍然影响后来者，一条路径让过去的移动影响下一次移动，一种叫声让危险在危险到来之前被传递。任何能让不在场之物继续影响行动的结构，都已经具有表示的影子。存在不再只在现场发生，它开始留下可被其他生命读取的痕迹。

人类把这种替代能力推进到新的层级。声音、手势、图案、刻痕、壁画，让事物第一次以可保存、可传播、可重排的形式出现。一头野牛不只在草原上存在，也开始在洞穴墙壁上存在；一次狩猎不只发生在现场，也开始在讲述中延续；一条河流不只流经土地，也开始进入方向、记忆和传说。神话、文字、账本、地图、法律和数字进一步扩大了这种能力：债务、土地、税赋、边界、王权和法律都可以离开现场，进入符号、记录和制度。一个命令可以跨越距离，一个账目可以跨越时间，一个法律可以在制定者死后继续约束行动。

\subhead{两种重量}

这时，世界中出现了一类新的存在。债务不是石头，法律不是身体，货币不是粮食，国家不是山川，但它们并不因此虚无。它们存在于记录、承认、执行、预期和责任之中\cite{searle1995construction}。它们没有自然物的直接质量，却能组织岩石、粮食、道路、矿山、军队和人的一生。这里必须区分两种重量：一种是物理重量，来自质量、体积、材料和能量；另一种是\term{本体论重量}，来自确立、维持、执行、回滚、外溢和承责的代价。社会性存在物理上可以极轻，却可能在现实中极重；它们轻于质量，却重于结算。

更准确地说，存在的重量不是一个单一刻度，而是由几种承付关系构成：确立、维持、回滚、外溢和承责。物质结构重在材料与能量，生命过程重在代谢、复制和不可逆扩散，社会结构重在承认、执行和信任，主体结构重在后果能否回流到一个持续边界中。所谓轻与重，不是价值高低，而是结算条件不同。把这几项支出加总，
\begin{equation}
W=w_{1}+w_{2}+w_{3}+w_{4}+w_{5},
\qquad
\frac{\partial W}{\partial w_{3}}>0,\quad
\frac{\partial W}{\partial w_{2}}>0
\label{eq:ontic-weight}
\end{equation}
式中 $W$ 是存在物的重量，不作绝对测量，只用来在同一套口径下比较大小，因此无量纲或随所选口径带量纲；$w_{1}$ 到 $w_{5}$ 依次是确立、维持、回滚、外溢、承责五项承付关系各自结算出的代价，五项必须折到同一口径（工时、能量、货币或注意力，任选其一）才能相加；$w_{3}$ 对应回滚，$w_{2}$ 对应维持。两个偏导说的是只增不减的两条方向：一个东西越难撤回、越难清除后果，$w_{3}$ 越大，$W$ 也就越大；一个东西越需要持续投入才能不退回噪声，$w_{2}$ 越大，$W$ 同样越大。五项在不同层次上的权重并不相同，这才是轻重的真正来源：同样一笔结算，落在物质结构上主要由 $w_{1}$ 与 $w_{2}$ 承担，落在主体结构上则主要由 $w_{3}$ 与 $w_{5}$ 承担。

科学、工程和计算把这种分层推到新的深度。一座桥可以先存在于图纸中，一座城可以先存在于规划中，一个分子可以先存在于模型中，一个气候系统可以先存在于仿真中。图纸不是桥，公式不是机器，模型不是现实，但它们让人类能够在物质动工之前排列可能性。计算机的特殊性在于，它不只保存表示，还能运算表示：文字、图像、声音、地图、账本、公式和图纸，都可以转化为比特；模型可以被修改，参数可以被调整，未来可以被推演，现实可以在行动之前被转化为候选空间。这类存在既不属于材料，也不属于纯粹虚构，它们的重量要按另一套条件结算；机器智能把这类存在的数量推向新的量级。

\subhead{生成物的层级}

层级说明的是不同存在的重量，机器智能改变的则是生成。它把这条道路推到了新的强度。过去，生成一个世界需要神话，生成一座城市需要图纸，生成一个人物需要小说，生成一个分子需要实验室，生成一个制度需要漫长的政治想象。现在，文字、图像、声音、代码、角色、方案、药物候选、城市规划和未来情景，都可以在模型中快速出现。表示不再只是保存现实，也不只是预演现实，而开始大规模生成可能的现实。但这并不是一个全新的本体论问题。人类很早就知道，图纸不是桥，地图不是土地，法律文本不是社会秩序，货币符号不是粮食本身，小说人物不是现实人物。机器智能的特殊性，不在于第一次提出“生成物是否存在”，而在于它让表示性存在以远高于以往的规模和逼真度膨胀，把一个古老问题推到日常生活和技术系统的前台。

把生成物按层级摊开之后，追问的方向就变了。真正的问题不是：被生成的东西算不算存在。这个问题太粗。答案也太简单：它们存在，但以某种方式存在。

一个被生成的城市，作为图像和方案存在；它还没有土地、建筑、道路、居民、产权、维护和财政。一个被设计的分子，作为候选结构存在；它还没有合成、纯化、代谢、安全性、剂量和临床后果。一个仿真中成功的机器人动作，作为模型状态存在；它还没有摩擦、碰撞、磨损、误差和责任边界。一个语言中完整的人格，作为叙事对象存在；它还没有身体、时间连续性、后果回流和责任承担。它们不是不存在，而是存在于较轻的层级。

智能时代真正尖锐的问题，不是虚拟物是否存在，而是它存在于哪一层、经由哪条边进入现实、由谁承付后果。一个东西可以作为文本存在，作为图像存在，作为代码存在，作为仿真存在，作为候选存在，作为物理结构存在，作为生命过程存在，作为社会制度存在，作为责任对象存在。不同层级的存在，承受不同约束，要求不同维持条件，具有不同回滚难度。表示性存在能被说出、画出、生成、复制和修改；物质性存在要占据空间、承受负荷、消耗能量、维护结构；生命性存在要代谢、修复、适应、复制并反馈；社会性存在要被承认、执行、信任和追责；主体性存在则必须在时间中维持边界，并把行动后果重新带回自身。层级之间如何通行，由此成为比层级本身更紧的问题。

\subhead{重量与可达性的解耦}

技术黏滞与本体论在这里相接。技术黏滞不只是技术执行慢，不只是组织保守，也不只是资源不足。它更深的含义是：轻存在向重存在转换时，必须穿过材料、自由能、时间、结构、验证和责任。表示空间中的城市、分子、动作、制度和人格，可以快速生成；现实空间中的城市、药物、能力、制度和主体，必须承受世界的结算。机器智能降低的是候选出现的成本，不是目标层级的存在重量。

因此，存在结构不只是一座阶梯。阶梯说明了不同层级的重量：表示较轻，物质更重，生命更重，社会与主体承担更多结算条件。但智能暴露出另一个更重要的事实：层级之间不是只能逐级通过，它们之间还有\term{转换边}。有些边很重，有些边很轻；有些边过去几乎不存在，现在被技术打开；有些边看似低阻，实际把代价转移到别处。一条转换边之所以显得低阻，往往不是因为代价消失了，而是因为代价被延迟、外包，或转移到当下尚未显明的承受结构中。

层级与边分开之后，一个系统的位置就可以被确定。一个智能系统的本体论身份，不由它的技术架构决定，而由它的转换签名决定：它从哪一层生成，向哪一层结算，经过哪条边，谁承付维持代价、回滚代价和责任后果。生成式智能主要打开表示到社会的低阻边，文本、图像、角色和模型输出可以直接进入注意力、信任、舆论、组织流程和群体判断。具身智能打开表示到物质行动的执行边，模型策略一旦接入传感器、控制器和执行器，就必须面对摩擦、惯性、碰撞、能耗、误差和安全边界。还有一条更重的边，是表示进入生命过程：序列、结构和调控方案一旦被生命系统稳定读取、表达、复制，并进入遗传或生态传播，就不再只是候选，而进入代谢、发育、繁殖、演化和生态后果。

机器智能的本体论新意，不是制造了新楼层，而是改变了存在转换图的连边结构。更准确地说，\key{它让重量和可达性解耦。}旧世界里，重存在通常难以抵达：桥梁、城市、药物、制度、身份、责任，都需要材料、时间、组织、执行和承认，越重越慢，越重越难到达。机器智能打开的新通道，使某些轻载体能够低成本抵达高后果层。一个生成物自身可以很轻，但一旦改变人的选择、关系、预期和责任分布，它就在社会层变重；一个策略本身可以很轻，但一旦驱动车辆、机械臂或无人机，它就进入物质结算；一个序列本身可以很轻，但一旦被生命过程读取并复制，它就进入更深的不可逆性。

\subhead{重量的归属}

这不是轻存在逃过了结算。结算只是发生在别处：发生在人的注意力中，发生在平台规则中，发生在组织决策中，发生在信任结构中，发生在物体受力中，发生在生命过程里，发生在后来很难撤回的行为路径中。一个生成物自身可以很轻，但它调动的反应可以很重；它的载体可以被删除，但它造成的后果未必能回滚。因此，存在重量不能只看创建代价，也不能只看载体能否删除。更重要的是维持代价、后果回滚代价，以及这些代价最终归属于何种承受结构。最危险的存在形态，恰恰可能是载体极轻、后果极重。

更进一步，重量不是无主属性。物理重量的承受者是物理本身，岩石的质量不依赖人的承认，桥梁的承重和疲劳也不会因为无人相信而消失。但社会重量不同。债务重，是因为有人承认它、追索它、执行它；身份重，是因为制度和他人承认它；货币重，是因为交换网络承认它；机器角色重，是因为用户投入注意力、情绪、信任和时间。越往社会层和主体层走，重量越不是内在属性，而是关系量。同一个机器角色，对深度依赖它的人极重；对不识别它的系统几乎为零。同一个模型输出，对把它写进决策流程的组织极重；对外部旁观者只是文本。判断一个生成物是否变重，不能只看它的创建代价，也不能只看平台运行代价，还要看它的维持代价由谁承付。

智能还带来了另一种变化：表示开始形成自循环。过去，表示大多来自现实。地图来自土地，图纸指向工程，账本对应交易，照片来自现场，模型需要经验校准。今天，生成内容进入数据，数据训练模型\cite{shumailov2024ai}，模型继续生成内容，生成内容再进入人的注意力、平台分发和下一轮信息环境。表示不再只是现实的影子，而可能形成自我复制的回路。这不是虚拟替代现实，而是分叉：表示层可以在现实校准不足的情况下自我繁殖，同时又借助平台、关系和信任进入社会层，产生真实后果。最危险的形态，不是一个虚拟世界完全取代现实，而是一个没有充分现实校准的表示环境，持续影响人的判断、关系和行动。

在生命过程里，类似逻辑更重。一个进入生命过程的候选结构，如果被稳定读取、复制、遗传、扩散和选择，就可能不再依赖原始设计者持续维持。它开始沿着生命自己的机制传播。此时，智能生成的不是普通表示，而是进入自我复制过程的入口。这不是说智能可以任意创造生命。生命空间受到代谢、发育、稳定性、繁殖、生态和演化的强约束，可组合不等于可存活，可设计不等于可延续。真正的重量，不在序列被写出，而在它是否能被生命过程读取、维持、复制，并产生难以回滚的后果。

\subhead{本体论作为操作}

从这个角度看，本体论不是智能之外的抽象哲学，而是智能本身的底层操作。智能并不是在一个已经切分完成的世界里活动。它首先要把连续现实分解成可识别、可记忆、可预测、可操作的单位。生命把世界切分为有利与有害，动物把世界切分为猎物、威胁、路径和同伴，人类把世界切分为名称、工具、制度和责任，科学把世界切分为变量、模型、尺度和系统，计算机把世界切分为数据结构、状态空间和算法对象。机器智能则把这种切分推进到大规模生成、快速转换和跨层结算。

但这里隐藏着更深的裂缝。人类的本体论习惯于对象、边界、身份和责任。我们用“一个人”“一件物”“一个组织”“一个行为”“一个后果”来切分世界。可是机器智能的内部表示不完全以人类熟悉的离散对象为单位，它在很大程度上更接近连续空间中的方向、相似性、分布和可插值区域。它生成的对象，是人类界面上的对象；它内部处理的，未必是人类意义上的对象。人类需要对象，因为行动、责任、法律和记忆都依赖边界；机器智能可以在连续表示空间中处理相似性，却把结果输出成我们能理解的对象。这个差异暂时被界面掩盖了：我们看到的是图像、文本、角色、方案和回答，以为它在使用同样的本体论。但它内部执行的切分，未必与人类用来承担责任的切分一致。

因此，机器智能真正打开的问题，不是它生成了什么，而是它把多少轻载体送入了高后果层级，却没有同时生成相应的承受边界。一个模型可以影响公共判断，一个算法可以改变价格、排序、流量和机会，一个自动代理可以下单、筛选、推荐、拒绝、调度和执行，一个机器角色可以成为人的情感对象和决策参照。它们在社会层产生后果，却未必在主体层承受后果；它们可以参与行动，却没有身体，可以产生后果，却没有生命史，可以影响他人，却没有内在承受，可以被依赖，却没有稳定责任边界。

这些既产生后果又不受后果约束的存在物，撞上了社会的一条旧规则。传统社会默认一个基本等式：行动者应当也是承责者。自然人如此；公司、组织和国家虽然不是自然身体，但也通过法律、制度、财产、代表和责任链建立了责任结构，能够被起诉、惩罚、约束、解散、赔偿和追责\cite{dewey1926historic}。机器智能让这个等式松动了。开发者、部署者、平台、用户、组织、模型本身都可能参与后果链条，但责任在它们之间分散、转移、稀释，最后变成无人完整承受的残差。更深地说，机器智能的可回滚性，是它能力的一部分，也是它无法自然承责的原因。它可以快速生成、复制、替换和重置，因此能以极低成本参与判断和执行；但承责需要相反的东西：一个会被后果改变、不能随意逃逸、能够在时间中持续的边界。

这就是本体论通向主体论的地方。机器智能制造的不是一个简单的“像不像人”的问题，而是一类有社会后果、物质后果甚至生命后果，却未必有主体边界的存在物。它们不是虚无，也不是完整主体；不是普通工具，也不是传统制度法人；不是生命，却能进入关系；不是承责者，却能改变后果。智能时代最危险的本体论错误，不是相信虚拟物“存在”，而是忘记它落在哪一个存在层级上、靠哪一条边维持、又由谁承担它维持不住时的代价。

从宇宙太初到机器智能，存在的历史可以被看成一条漫长的分层之路，也是一条不断开辟新通道的路径。这条道路没有错。没有表示，现实只能被动发生；有了表示，现实可以被预演、计划和改造。智能之所以强大，正是因为它能让世界在行动之前先变成可处理的对象。

但这条道路也有边界。表示越轻，越容易忘记现实的重量；生成越快，越容易低估转换的代价；智能越能制造可能性，越需要区分可能性所在的层级，识别通向现实的路径，并追问后果由谁承受。世界可以被表示得越来越快，但现实仍然要求存在被材料承载，被时间展开，被生命反馈，被社会承认，被责任接住。生成不是终点。存在要被维持，后果要被回滚，责任要有承受者。

到这里，问题不再只是“什么存在”，而是：代价归属于何处，谁在维持，谁在行动，谁来承受。而在追问代价归属之前，还有一个更早的问题需要回答。被切分、保存、传播和生成的这些表示，究竟凭什么能够被依赖。

\subsection{知识的凭据：候选之后，什么能够被依赖}

表示被说出、保存和依赖的规模越大，它凭什么可以被依赖，问题就越硬。存在可以按层级摊开，一种东西可以很轻，却仍然抵达很重的后果层；在这一步之前，还有一件更早的事需要交代：在后果真正发生之前，行动所依据的判断凭什么站得住。

\subhead{校准先于命题}

知识不是世界本身，也不是表示本身\cite{ogden1923meaning}。世界并不因为被理解才运行，表示也不因为被说出就成为知识。对生命来说，最早的“知道”不是命题，而是校准：什么能接近，什么会破坏；什么能维持边界，什么会带来损伤。错误不是逻辑问题，而是生存问题。把毒物当成营养，代价不是观点被反驳，而是结构被破坏。由此看，知识从一开始就不是纯粹观看，而是在世界约束下形成可依赖差异的过程。

这里讨论的\term{凭据}，首先指向那些会被依赖、会指导行动、会进入现实后果链的知识。纯数学和形式逻辑是一个极限情形：它们的凭据主要来自证明、结构一致性和公理边界\cite{godel1931}，而不是经验接触；但只要它们被用于工程、计算、金融、治理或科学建模，就会重新进入现实后果链。换言之，凭据的形式可以不同，但任何被依赖的判断，都必须说明自己在何种结构内成立、何时可以迁移、何处必须停止。人类让校准离开身体，进入语言、文字、图像、数字和制度记录。经验因此可以传播、积累和复用，也可以伪造、漂移和固化偏见。真正重要的不是人类会说，而是人类逐渐发明了让说法接受约束的方式。测量让感觉进入尺度，实验让判断暴露在结果面前，复现让一个说法不能只依赖某个权威，数学和仪器把直觉无法稳定分辨的差异压入可比较、可推演、可检验的结构。科学的特殊之处，不在于它让人类说出了更多东西，而在于它发明了一套让世界有机会拒绝人类说法的机制\cite{popper1934logik}。

\subhead{凭据的两层结构}

从这个意义上说，知识不是一种更漂亮的表示，而是一种获得凭据的表示\cite{gettier1963justified}。凭据不是由流畅性、完整性、自洽叙述或权威语气授予的，而来自两层结构：第一层是形成层，即接触、校准、复核；第二层是兑现层，即边界和后果。接触让候选暴露于约束，校准让误差返回判断结构，复核让判断离开单点输出；边界说明它在哪里成立、在哪里失效，后果说明它被依赖之后会改变什么、由谁承受什么代价。两层缺一不可：只有形成过程而没有边界与后果声明的候选，无法说明自己何时失效；只有后果要求而没有接触与复核的判断，只能靠权威维持。

凭据如何形成？不同介质会改变凭据形式，但不改变基本结构。接触让候选暴露于对象、材料、身体、环境、行为或结构约束；校准让误差返回判断尺度，使下一次判断因此改变；复核让判断离开单点输出，进入不同主体、工具和场景的公共检验。成熟知识不是无边界的自信，而是知道自己在哪里成立、在哪里失效；一个判断一旦被依赖，就会进入行动，改变路径、资源、身体、组织和未来。没有后果要求的候选可以作为想象、草案、启发或游戏；一旦进入身体、工程、交通、金融、治理或生命过程，就不能只凭形式完整获得通行。

凭据的形成与兑现不排斥机器，它们只是规定了机器能做什么。机器当然可以参与验证。它可以运行测试，比较数据，发现异常，辅助复核，搜索反例，执行仿真，检查一致性，甚至在某些封闭环境中完成高度自动化的确证流程。问题不在于机器能不能参与验证，而在于验证的最后代价由现实接触设定，验证的充分性标准则由后果结构决定。合成要接触材料，治疗要接触身体，交通要接触道路，工程要接触负荷，生命干预要接触代谢与生态，制度设计要接触人群行为。机器可以加速某些环节，却不能自动决定“多少凭据才算足够”，因为这个问题最终要回到边界、后果、风险和承受者。

因此，智能时代的知识问题，不是机器是否“真的理解”。这个问题仍然绕在人类心智的旧词汇里。更稳的问题是：一个输出是否被置于足够的接触、校准、复核、边界和后果约束之中；它是否知道自己的适用范围，是否能够暴露来源，是否能够标明条件，是否能够在失败时被追踪，是否能够被现实修正。知识不是一个孤立答案，而是一条可追溯的形成路径。这条路径的起点不在表示内部，而在表示之外的接触、工具、场景和制度。

\subhead{候选扩张与凭据带宽}

机器智能改变了这条历史的压力分布。从认识论结构看，它是文字发明以来规模最大的一次表示技术扩张。文字让记忆离开身体，人的认识位置从背诵转向理解、辨析和解释；印刷让复制脱离手抄，人的认识位置从占有文本转向筛选、校验和诠释；计算让形式推演和搜索脱离心算，人的认识位置转向模型边界、数据来源和结果验证。机器智能延续并放大了这条线：它降低的不是知识成立的成本，而是候选出现的成本。正因为候选可以被大规模、低成本地提出，人类在认识论中的位置并没有消失，而是从答案的生产者转向凭据结构的设计者、边界的判断者，以及可依赖性本身的形成者。

更具体地说，机器智能没有首先取消无知，而是极大扩张了候选空间。这里的候选不只是文本回答或屏幕图像，也可以是解释、代码、控制策略、实验方案、分子结构、蛋白构象、交通路径、诊断判断、金融操作、制度草案，乃至进入生命过程的序列和调控方案。智能制造的不是已经成立的知识，而是有待检验的可能路径。介质不同，约束不同，验证方式不同，但认识论问题是同一个：这些候选凭什么被依赖？

\key{智能扩张候选，验证形成凭据。}

这句判断的两侧，现在可以分别结算了。这是智能时代认识论的核心，也是上一节“重量与可达性解耦”在认识论上的对应面。前面说的是轻载体如何低成本抵达高后果层；这里讨论的是候选如何低成本出现，却未必携带相应凭据。存在论上，问题是轻东西如何产生重后果；认识论上，问题是快候选如何获得慢凭据。机器智能真正改变的，不是让生成自动等于存在，也不是让生成自动等于知识，而是把可能性出现的速度同现实承付的速度分离开来。过去，稀缺常常在信息获得；现在，许多领域的表述能力、组合能力和搜索能力被急剧放大，稀缺开始转移：不再是有没有说法，而是哪一个说法经过了足够约束；不再是有没有路径，而是哪一条路径可以被依赖。更准确地说，机器智能使\term{候选通量}与\term{凭据带宽}发生解耦：候选出现的速度被大幅提高，而接触、校准、复核、边界和后果仍然受材料、身体、时间、制度和责任限制。由此产生两类相反的失败：在存在层面，轻载体可能已经抵达高后果层，却没有明确承受者；在认识层面，候选尚未穿过现实，却已经冒领了知识外形。

把这两个速率写进同一个流量关系，

\begin{equation}
\frac{dQ}{dt}=\lambda_{\text{候选}}-\lambda_{\text{凭据}}
\end{equation}

式中，$Q$ 是尚未穿过接触、校准、复核、边界声明和后果检验的候选存量，单位为个；$t$ 是时间，单位为年；$\lambda_{\text{候选}}$ 是候选出现的速率，$\lambda_{\text{凭据}}$ 是凭据形成的速率，二者单位均为个每年，故右侧量纲回到个每年。这个式子的含义是：解耦不需要额外的机制，它只是两个速率不再相等。当候选出现的速率高于凭据形成的速率时，存量按两者之差单调增长，而增长的存量不会自己生成凭据；这就是“不缺说法、缺的是可以被依赖的说法”的结构形式。适用条件是两侧必须折算到同一粒度，候选与凭据只有在能按同一种单位计数时，差值才有意义；而且它只给出趋势方向，不给出长期预测，因为存量增长会反过来改变筛选标准、注意力分配和制度安排，$\lambda_{\text{凭据}}$ 本身随存量变化。它刻画的是结构，不提供任何实测数值。

\subhead{形式外壳与认识论黏滞}

可依赖性因此成为新的核心稀缺。一个回答可以很快，一个判断不能因此可靠；一个解释可以圆满，一个机制不能因此成立；一个动作可以在仿真中成功，一个真实系统不能因此安全；一个序列可以被写出，一个生命过程不能因此可控。机器智能让候选获得了知识的外形，却没有自动给予它被依赖的凭据。危险也正在这里：错误不再必然粗糙，无知不再必然沉默，虚假不再必然破绽明显。它可以流畅、礼貌、系统、专业，并且在表面上完成所有表达动作。因此，这不是简单的“幻觉”问题；更深的变化是候选的生产机制与确证机制发生了分离，形式的完成速度超过现实的校准速度，可能性的展开速度超过凭据的形成速度。

这种变化会改变人类对知识的直觉。人类长期把流畅表达、概念组织、语气稳定、格式规范和逻辑连贯当作可信信号，因为在过去，这些信号通常需要训练、经验和制度背景才能产生。机器智能打破了这种关联：它保留了可信信号的外壳，却降低了外壳的生成门槛。于是，认识论的任务不再只是区分有无信息，而是区分候选是否获得凭据。在这个意义上，技术黏滞也可以被理解为认识论黏滞。不是世界拒绝智能，而是世界要求判断付出接触成本；不是现实落后于生成，而是现实要求任何候选在被依赖之前先承受约束。生成可以在表示空间中快速展开，验证必须接触世界；生成可以并行制造大量候选，验证必须面对材料、时间、实验、身体、场景、制度和责任。生成常常是轻的，验证却会变重，因为验证要让候选穿过现实的阻力。

表示自循环的本体论后果已经讨论过：表示可以在现实校准不足的情况下自我繁殖，并通过平台、关系和信任进入社会层，形成一种轻载体、重后果的存在分叉。这里要看的是它的认识论后果。一个认识系统如果主要从自身的回声中获得材料，就会逐渐削弱接触比例；它可能越来越擅长内部一致，越来越不擅长承受外部反驳。仿真、指标、模型、平台反馈和组织报告，都可能形成类似的自我喂养回路\cite{cinelli2021echo}。系统越是用内部生成物替代外部接触，就越可能形成高度一致的候选世界，却不再知道现实从哪里打断它。知识不是被现实支持的漂亮表示，而是保留了反驳接口、并能把反驳转化为校准的候选。

\subhead{验证的分配}

候选爆炸之后，真正困难的还不只是验证，而是验证如何被分配。接触世界需要材料、实验、时间、场景、组织、注意力和责任，不可能平均分配给所有候选。因此，人类的新认识角色不是逐一验证所有输出，而是在候选洪流前设计验证路径：哪些候选因为后果不可逆而必须优先检验，哪些候选因为越出边界太远而必须暂缓依赖，哪些错误因为暴露延迟而需要更高等级的复核，哪些判断只能停留在启发层，哪些可以进入行动准入。验证不是单纯的技术动作，而是一种稀缺接触预算的分配。这也是人类知识制度的价值所在。图书馆、实验室、大学、标准体系、期刊、审计、临床试验、工程认证、事故调查、法院证据、同行批评，并不是知识之外的附属物，而是凭据形成和验证分配的结构\cite{merton1942normative}。它们缓慢、昂贵、烦琐，有时保守，有时僵化，但它们承担一个基本功能：防止任何流畅的表示、漂亮的模型或成功的仿真绕过接触，直接冒充可依赖的判断。

机器智能不会让这些结构消失，相反会使它们更重要。当候选说法稀少时，人类可以把大量注意力放在生成上；当候选空间膨胀时，人类必须把更多注意力转向筛选、验证、追溯、边界判断和验证优先级。未来真正稀缺的能力，不只是提出答案，而是为答案建立凭据；不只是生成方案，而是知道方案需要通过哪些现实约束；不只是给出结论，而是判断哪些结论值得消耗昂贵的现实接触。验证预算有限，这决定了认识论的主要工作从生产转向排序。

所以，智能时代的认识论不能只问真假。真假当然重要，但更第一性的问题是：这个候选经过了什么约束？它的凭据从哪里来？它能被谁复核？它的失败如何暴露？它的边界如何声明？它一旦被使用，会进入什么后果链？

如果这些问题没有答案，那么它最多是候选，而不是知识。候选并不低贱，人类所有知识都从候选开始；但候选不是终点。认识论的核心任务，就是不要让候选未经凭据化就伪装成知识，不要让形式的完整性冒充现实的通过。

这也意味着，人类在智能时代的认识角色不会消失，而会改变。人不再只是答案的生产者，而越来越成为验证路径的设计者、边界的判断者、接触预算的分配者和可依赖性本身的形成者。机器可以扩大候选，也可以执行大量验证环节；但哪些候选值得接触现实、哪些判断可以被依赖、多少凭据足以行动、哪些后果必须被优先看见，这些问题最终会回到主体、价值和承受者。只有把机器智能放在正确位置上，它才真正强大：把候选当成知识，会制造虚假的确定性；把候选放回候选，则能极大扩展探索空间。

这就是问题核心：候选之后，什么能够被依赖？答案不在于它是否最流畅、最像专家、最符合期待、最能安慰人，也不在于它是否最快给出结果。这些只是表达层面的外观，不能构成知识的凭据。真正能够被依赖的，是那些经过接触、校准、复核、边界声明和后果检验的判断、结构与路径。知识不是语言的胜利，而是候选在现实约束中存活下来的形式。机器智能之后，人类面对的不是知识自动丰收，而是候选空前丰盛。真正稀缺的，是让候选穿过现实、形成凭据、成为可依赖判断的能力。从这里开始，问题不再只是“什么被生成”，而是：什么能够被验证，什么可以被依赖，什么可以获得进入行动的准入。

\subsection{行动的准入：表示何时可以跨过转换边}

候选之后，需要凭据；凭据之后，仍然不能直接行动。知识与行动之间不是一条连续斜坡，而是一条转换边。知识使一个判断获得可依赖性，行动则使这个判断进入现实过程。前者仍然处在表示关系之中，后者已经开始改变世界状态。一个模型可以有统计优势，一个方案可以通过局部测试，一个控制策略可以在仿真中稳定，一个制度设计可以在理论上自洽，但这些都还没有回答更重的问题：它是否可以消耗资源、占用时间、触碰身体、调动组织、改变制度，并使残余后果落到具体承付结构上。

\subhead{转换边}

前面两节分别讨论了存在的重量和知识的凭据：轻载体可以抵达高后果层，重量与可达性被机器智能重新拉开；候选可以高速生成，但凭据不能同速形成，候选通量与凭据带宽发生脱钩。到这里要处理的是第三个断裂：表示可以获得现实触发力，却不自动携带后果承付结构。核心问题不是“能不能输出”，也不只是“能不能被相信”，而是一个表示在什么条件下可以触发现实状态转移。

\key{触发力与承付结构的解耦，是行动论在机器智能时代遇到的新问题。}

承付结构在传统世界里由行动者本人提供。行动通常与一个有限主体捆在一起。一个人做出行动，不只是发出一个指令或表达一个愿望，而是把身体、身份、声誉、责任、时间和未来押入某条路径。行动者不是纯粹的触发点，也是后果回流的边界。失败会改变他，损失会落到他身上，责任会追溯到他，未来义务会继续缠住他。正因为行动者自身被押入后果，行动才不是轻飘的表示，而是有重量的选路\cite{arendt1958human}。

机器智能改变的不是这个结构的意义，而是它的耦合方式。模型可以生成指令、调用工具、修改文件、发送信息、执行代码、访问数据库、触发审批、改变排序、调度设备。它可以通过权限、接口和工作流触发现实变化，但这些变化未必回流到模型本身。模型可以拥有触发力，却没有身体损伤；可以参与路径选择，却没有生活世界；可以放大后果，却没有资产边界、声誉连续性和惩罚敏感性。触发发生了，承付却不在同一边界内完成。

\subhead{行动论的边界}

这就是为什么传统行动论既重要，又不能主导本章。传统行动论长期追问：什么使一个事件成为行动？行动是否由意图、理由、欲望、信念或计划组织\cite{davidson1963actions}？一个行为为什么可以归属于某个行动者？这些问题背后真正处理的，是行动归属与后果回流的耦合。意图、理由和计划并不只是内心事实，它们把“为什么选这条路”“谁在选路”“谁延续这个选择”连接起来。一个有计划的行动，是主体用未来约束现在，也是把后果提前纳入自身连续性\cite{anscombe1957intention}。

AI agent 使这个框架遇到边界。问题不必首先落在“AI 是否拥有真正意图”上。即使没有传统意义上的完整意图，系统仍然能够触发现实状态转移。它可以移动资金、驱动设备、触发流程、传播信息、改变排序、重排机会。于是行动论的问题发生偏移：不再只是“这个行为是否属于某个主体”，而是“某个表示通过什么接口进入现实，获得了多大作用半径，并由什么结构承付它的残余后果”。

这也是近年 AI agent 使风险形态发生变化的原因。普通生成模型首先带来内容风险：错误、幻觉、偏见、误导、伪造。agent 则把问题推进到行动边界：权限、身份、日志、回滚、停机、监控、持续运行和责任范围\cite{openai2025chatgptagentsystemcard}。过去，模型输出离行动很远，中间隔着阅读、理解、转写、审批、执行、复核等多重摩擦；现在，读、写、查、调、改、发起、调用和执行可以被连成同一条工作流。模型不只是说出可能性，而是在某些接口中获得触发力。

这说明，agent 的实质不只是模型更聪明，而是表示跨过转换边的摩擦被降低了。过去，组织摩擦在无形中延缓了错误进入现实；今天，自动化把这种摩擦压缩为权限调用。过去，一个建议要经过人类转译才可能行动；今天，建议、调用和执行可以在同一系统中连续发生。风险焦点因此从“表示是否正确”转向“表示如何获得触发力”。行动准入也因此不再是伦理附录，而是机器智能进入现实系统的第一边界。

\subhead{知识与行动的分界}

要把这一点说清楚，必须区分知识与行动。知识降低不确定性，行动花费不可逆性。验证消耗的是接触预算，行动消耗的是不可逆预算。一个判断越可靠，越有资格被依赖；但被依赖仍不等于可以执行。因为判断仍然属于表示空间，行动则进入现实历程。表示空间允许多个候选并存，现实系统只能沿具体路径演化。一个方案进入执行，其他方案就被延迟、排除或失去条件；一个政策进入实施，社会激励就开始重排；一个医疗方案进入身体，代谢和免疫过程就已经启动。知识管的是“能否信”，行动管的是“是否花掉一段不可逆”。

行动不是知识的应用，而是不完备知识对现实的暴露。任何行动都不可能等到未来完全展开之后才发生。再充分的凭据也只能降低不确定性，不能把未来全部提前完成。行动发生时，总有一部分残余不确定性从模型、文本、计划或仿真内部转移到现实系统之中。这个转移不是修辞意义上的风险，而是结构意义上的状态改变：未知后果开始由身体、设备、市场、组织、制度或生态过程承载。

因此，行动的重量不来自动作本身，而来自它对可达空间的处分。表示空间可以容纳很多尚未实现的可能；行动则把其中一条路径实际化，同时消耗资源、占用时间、改变结构，并关闭一部分未来可能性。一个医疗建议变成用药，身体就进入特定代谢路径；一个交易策略变成自动执行，风险就进入市场结构；一个招聘模型变成筛选规则，某些人的机会空间就被提前收缩；一个调度策略进入电网、交通或工厂，物理系统就必须沿它给出的节奏响应。即使局部后果可以回滚，时间、机会、信任、损伤和路径依赖也不能完全恢复为行动之前的状态。

\subhead{行动与可达空间}

这使讨论重新回到结构与可达空间。第二章讨论的结构记忆已经说明撤回不等于恢复；同一章的可达空间分析已经说明系统并非在无限可能中自由移动，而是在历史路径、鲁棒性和结构约束中保留有限可塑性。行动正是可达空间被花掉的动作。每一次不可逆行动，都在某种程度上改变未来可达域。机器智能让选路、触发和执行变得更轻，因此也让“花掉不可逆性”变得更容易。如果没有准入，系统会以很高速度消耗自己的可塑性。

这也把行动准入接回 RTG 模型。行动一旦发生，R 不再只是资源估算，而是资源被实际占用、锁定或消耗；T 不再只是等待或验证周期，而是不可逆历程；G 不再只是结构关系图，而是被实际牵动的反馈、耦合和责任链。行动是 RTG 从分析坐标变成现实结算的时刻。它把可能路径转换为现实路径，也把模型残差转换为系统残差。

把行动放回 RTG 之后，还要区分它穿过的是哪一种边，因为不同转换边具有不同重量。表示进入社会判断，首先改变注意力、信任和预期。错误图像、伪造引用、看似专业的解释、被自动分发的结论，可以进入人的信念结构，继而改变组织和公共讨论。这里的后果不一定表现为直接物理伤害，而是判断环境被污染。

表示进入物理执行，后果更直接。控制策略一旦接入车辆、机械臂、电网、无人机、工厂或医疗设备，就不再只是概率分布，而是力、热、速度、位移、碰撞、磨损和故障。表示进入生命过程，转换边更重。生命不是被动载体，而是会读取、代谢、放大、修复、抵抗、代偿和演化的过程。表示进入制度决策，后果最容易被低估。一个评分模型进入招聘，一个风险模型进入信贷，一个预测系统进入警务，一个分配算法进入教育或医疗，都不是单纯提供建议，而是在重排人的可达空间\cite{oneil2016weapons}。

\subhead{承付的重建与持续准入}

因此，行动准入不能被理解为“多加一层安全要求”。它的功能，是在触发力与承付结构被切开的地方，人工重新缝合它们。旧世界里，主体常常免费提供这道闸门：行动者把自己押入后果，行动因此天然带有重量。机器智能使触发可以脱离这种押入，准入就必须把承付结构重新挂回行动点。没有这个重建，行动就会成为无承付的触发。

把承付挂回行动点有几种做法。可验证、可停止、可回滚、可追溯、可承责，不是平铺的治理清单，而是重建承付结构的不同方式。没有验证，残差不可识别；没有停止，残差不可中断；没有回滚，残差不可收束；没有追溯，残差不可定位；没有承责，残差没有回流边界。它们共同回答的不是“系统是否合规”，而是“未消除的不确定性进入现实后，是否有承载、收束和回流结构”。

准入不是给知识盖章，而是在触发之前重建承付。

据此可以说，“人在回路中”本身不够。关键不在于人是否形式上出现，而在于人是否拥有真实的承付位置。一个只能点击确认的人，不是准入者，只是自动化流程中的责任垫片。一个看不见模型边界、看不见后果链、无法改变执行路径、又必须在极短时间内确认的人，不是在监督系统，而是在替系统承担名义责任。有效监督要求人能够看见边界、理解残差、改变路径、中止过程，并实际处在后果回流链条中\cite{santoni2018meaningful}。

更困难的是，转换边常常不是一次性跨过的。很多系统开始只是辅助，后来成为推荐；开始只是推荐，后来成为默认；开始只是默认，后来成为自动；开始只是自动，后来变成无人真正理解和质疑的基础设施。表示不是突然变成行动，而是被界面、流程、绩效压力和组织惯性逐步推近现实。最后出现的不是一次越界，而是一条已经固化的路径。

对于持续运行的 agent，行动准入也不能只发生在部署之前。能够调用工具、读取反馈并调整策略的系统，不是在执行一次动作，而是在形成一条持续作用的因果管道。准入必须规定作用半径、权限边界、更新频率、监测方式、停机条件和责任回路。一次测试通过，只能说明系统在某些条件下可以开始，不能说明它在所有未来反馈中都拥有不受限制的行动权。持续行动需要持续准入。在这个意义上，准入评价的不是一个表示是否精彩、连贯或有启发性，而是它获得触发力之后，是否具备相应的承付结构。若剩余不确定性找不到明确的承担者，表示仍可作为候选、解释、模拟或辅助存在，但不应获得现实执行权。行动的门槛，不是知识的完成，而是残差的可承付化。

到这里，前面的线索合在一处：轻载体如何抵达重后果，快候选如何获得慢凭据，表示获得现实触发力时又如何必须重建承付结构。否则，机器智能会生成更多候选，也会更快花掉不可逆性，更快压缩系统可达空间，并把残差外部化给那些没有准备承受它的身体、组织和制度。行动把残余不确定性推入现实，并使其落在某些位置上；但后果落在哪里，并不等于谁在承受。身体、市场、组织、平台、部署者、用户和受影响者都可能被迫承载后果，却未必形成清晰的回流边界。于是，行动准入之后，问题必然推进到下一层：什么样的边界，才有资格被称为后果的承受者？
